\documentclass[11pt, a4paper, oneside]{article}
\usepackage[ngerman]{babel}
\usepackage{worksheet}

\title{Webentwicklung: \hspace{10cm} Javascript}
\author{L. Bung}
\subject{SAE}
\class{E3FI}

\begin{document}
	\maketitle
	
	
	\begin{hint}{Was ist JavaScript?}
		JavaScript ist eine skriptbasierte Programmiersprache, die standardmäßig im Browser ausgeführt wird.
		Während HTML die Struktur und CSS das visuelle Design definieren, sorgt JavaScript für dynamisches Verhalten, Logik und Benutzerinteraktion.
	\end{hint}
	
	
	\singletask{Grundlegende Syntax und Variablen}
	
	JavaScript besitzt eine C-ähnliche Syntax, die sich schnell in bestehende Programmierkenntnisse integrieren lässt.
	Betrachten Sie den folgenden kurzen Codeausschnitt.
	
	\begin{lstlisting}[language=javascript]
const inventory = [
	{ id: 1, item: "Laptop", stock: 5 },
	{ id: 2, item: "Monitor", stock: 0 },
	{ id: 3, item: "Dockingstation", stock: 12 }
];

for (let i = 0; i < inventory.length; i++) {
	let curr = inventory[i];
	if (curr.stock === 0) {
		console.log("Artikel nachbestellen: " + curr.item);
	} else {
		console.log("Verfügbar: " + curr.item + " (" + curr.stock + ")");
	}
}
	\end{lstlisting}
	
	\begin{enumerate}[label=\alph*)]
		\item Erklären Sie den Unterschied zwischen den Schlüsselwörtern \texttt{const} und \texttt{let} bei der Deklaration von Variablen.
		\item Beschreiben Sie die Funktion des Codes und schreiben Sie den Javascript-Code in Python um.
		\item Recherchieren Sie die Syntax von Programmstrukturen, die Sie bereits aus Python kennen: Verzweigungen, while- und for-Schleifen, Listen...
	\end{enumerate}
	
	\lines[3.5cm]
	
	
	\partnertask{DOM-Manipulation}
	
	Um dynamisch mit einer Webseite zu interagieren, muss JavaScript über das Document Object Model (DOM) auf einzelne HTML-Elemente zugreifen können.
	Zwei grundlegende Methoden hierfür sind \texttt{document.getElementById()} und \texttt{document.getElementsByClassName()}.
	
	\begin{enumerate}[label=\alph*)]
		\item Erläutern Sie den Unterschied zwischen diesen beiden Methoden.
		\item Gegeben sei ein HTML-Absatz mit der ID \texttt{status-message}.
		Schreiben Sie eine JavaScript-Anweisung, die den Textinhalt dieses Elements auf den Wert ``Bereit'' ändert.
	\end{enumerate}
	
	\lines[4cm]
	
	
	\task{Eigenständige Entwicklung eines Codebeispiels}
	
	In dieser Aufgabe entwickeln Sie ein eigenes kleines Skript, das auf eine Benutzerinteraktion reagiert.
	
	\begin{enumerate}[label=\alph*)]
		\item Skizzieren Sie ein minimales HTML-Grundgerüst, das einen Button sowie ein auswertbares Ausgabeelement (z.\,B. ein \texttt{<div>} oder \texttt{<p>}) mit einer eindeutigen ID enthält.
		\item Schreiben Sie den passenden JavaScript-Code, der beim Klick auf den Button eine Aktion ausführt (beispielsweise das Umschalten eines Textes oder das Hochzählen eines numerischen Wertes).
	\end{enumerate}
	
	
	\task{Clientseitige Formularvalidierung}
	
	In der Webentwicklung ist es üblich, Benutzereingaben direkt im Browser (clientseitig) zu validieren, bevor sie an den Server geschickt werden, um unnötige Requests zu vermeiden.
	Gegeben ist ein einfaches Eingabefeld für ein Passwort sowie eine Fehlerausgabe.
	
	\begin{enumerate}[label=\alph*)]
		\item Skizzieren Sie die Logik für eine JavaScript-Funktion, die beim Absenden eines Formulars prüft, ob das eingegebene Passwort mindestens 8 Zeichen lang ist und mindestens eine Zahl enthält.
		\item Erläutern Sie, wie Sie mithilfe von JavaScript den Text der Fehlerausgabe dynamisch ein- oder ausblenden bzw. anpassen können, wenn die Validierung fehlschlägt.
	\end{enumerate}
	
	\lines[5cm]
	
	
	\bonustask{Spezialitäten und unerwartetes Verhalten in JavaScript}
	
	JavaScript ist historisch bedingt für einige sprachliche Eigenheiten bekannt.
	
	
	\begin{enumerate}[label=\alph*)]
		\item Analysieren Sie die folgenden Ausdrücke und erklären Sie kurz, was die Konsole ausgibt bzw. warum dieses Verhalten auftritt:
		\begin{enumerate}[label=\roman*)]
			\item \lstinline[language=javascript]{console.log("5" + 3);} im Vergleich zu \lstinline[language=javascript]{console.log("5" - 3);}
			\item \lstinline[language=javascript]{console.log(typeof null);}
			\item \lstinline[language=javascript]{console.log(0.1 + 0.2 === 0.3);}
		\end{enumerate}
		\item Finden Sie selbst weitere Eigenheiten und zeigen Sie diese Ihren MitschülerInnen.
	\end{enumerate}
	
	\lines[4cm]
	
	\begin{goallist}
		...Javascript-Code lesen und erklären. &&&\\
		\hline
		...Beispiele für sinnvolle und ungeeignete Nutzungen von Javascript nennen. &&&\\
		\hline
		...HTML-Elemente mithilfe von Javascript-Code gezielt modifizieren. &&&\\
		\hline
	\end{goallist}
\end{document}